% Chapter 01: Kit Anatomy
% Directory structure, manifest format, versioning contract

\chapter{Kit Anatomy}
\label{ch:kit-anatomy}

% ============================================================
\section{Directory Structure}
\label{sec:kit-directory}
% ============================================================

A Domain Kit is a directory containing a manifest, a compiler binary, schema
artifacts, validators, projections, and test fixtures. The layout is
standardized to enable tooling (e.g., \texttt{gert kit test}) to discover Kit
components without additional configuration.

\paragraph{Canonical structure:}

\begin{verbatim}
gert.compliance/               # Kit root directory
├── gert-kit.yaml              # Kit manifest (REQUIRED)
├── bin/
│   ├── gert-kit-compliance    # Compiler binary (REQUIRED)
│   ├── validate-owner         # Optional: individual validator binary
│   └── project-audit-log      # Optional: individual projection binary
├── schemas/
│   └── compliance-v1.json     # Kit-specific JSON Schema overlay
├── test/
│   ├── fixtures/              # Test fixture runbooks
│   │   ├── 01-basic-approval.yaml
│   │   ├── 02-nested-governance.yaml
│   │   └── 03-invalid-missing-owner.yaml
│   └── expected/              # Expected lowered output and validation results
│       ├── 01-basic-approval.lowered.yaml
│       ├── 02-nested-governance.lowered.yaml
│       └── 03-invalid-missing-owner.errors.json
└── README.md                  # Optional: Kit documentation
\end{verbatim}

\paragraph{Required components:}

\begin{description}
  \item[\texttt{gert-kit.yaml}] --- The Kit manifest. Declares identity,
    compatibility, schema contributions, compiler entry point, validators, and
    projections. See \S\ref{sec:kit-manifest}.

  \item[\texttt{bin/<compiler-binary>}] --- The Kit compiler. A standalone
    executable that reads Kit YAML on stdin and writes core
    \texttt{runbook/v2} YAML to stdout. Must be named in the manifest under
    \texttt{compiler.executable}.
\end{description}

\paragraph{Optional components:}

\begin{description}
  \item[\texttt{schemas/*.json}] --- JSON Schema overlays that define
    Kit-specific step types and field structures. Referenced from the manifest
    under \texttt{schema.json-schema}.

  \item[\texttt{bin/validate-*}] --- Individual validator binaries. If your
    validators are bundled into the compiler, this directory may be empty.

  \item[\texttt{bin/project-*}] --- Individual projection binaries. These
    consume trace event streams (NDJSON on stdin) and produce projection
    output (YAML or JSON on stdout).

  \item[\texttt{test/fixtures/}] --- Test runbooks authored with this Kit's
    vocabulary. Used by \texttt{gert kit test}.

  \item[\texttt{test/expected/}] --- Expected lowered output and validation
    results for test fixtures. Used to assert correctness of the Kit compiler.
\end{description}

% ============================================================
\section{The Kit Manifest}
\label{sec:kit-manifest}
% ============================================================

The \texttt{gert-kit.yaml} manifest is a YAML document with \texttt{apiVersion:
kit/v1}. It declares the Kit's identity, compatibility constraints, schema
contributions, compiler location, and optional validators and projections.

\paragraph{Complete example:}

\begin{minted}{yaml}
# gert-kit.yaml — Domain Kit Manifest
apiVersion: kit/v1

meta:
  name: gert.compliance           # Reverse-DNS namespaced identifier
  version: "1.2.0"                # Semantic version (MUST be semver)
  description: "Compliance and audit-focused runbook authoring kit"
  author: "Acme Corp <compliance@acme.example>"
  homepage: "https://github.com/acme/gert-kit-compliance"

compatibility:
  gert-core: ">=2.0.0 <3.0.0"    # Semver range for gert core compatibility

schema:
  apiVersion-suffix: "compliance/v1"  # Authors use: runbook/v2+compliance/v1
  json-schema: "./schemas/compliance-v1.json"

compiler:
  executable: "./bin/gert-kit-compliance"
  args: ["compile"]              # Optional args passed to compiler
  input-format: yaml
  output-format: yaml

step-types:                      # Domain-specific step types this Kit defines
  - name: approval-request
    lowers-to: [manual, tool]    # Core types this compiles into
    description: "Request approval from named approvers with timeout"
  - name: evidence-capture
    lowers-to: [manual]
    description: "Capture structured evidence with attestation"

validators:                      # Optional domain-specific validators
  - name: compliance/owner-required
    phase: parse                 # parse | plan
    description: "Every compliance runbook must declare an owner in meta.roles"
  - name: compliance/evidence-on-sensitive
    phase: plan
    description: "Steps marked sensitive must capture evidence"

projections:                     # Optional read models over trace data
  - name: compliance/audit-log
    description: "Audit log view with approver attestations and evidence links"
  - name: compliance/evidence-bundle
    description: "Compliance evidence bundle in JSON format"
\end{minted}

\paragraph{Required fields:}

\begin{description}
  \item[\texttt{apiVersion}] --- MUST be \texttt{kit/v1}.
  \item[\texttt{meta.name}] --- Reverse-DNS namespaced identifier (e.g.,
    \texttt{gert.compliance}, \texttt{com.acme.workflow}).
  \item[\texttt{meta.version}] --- Semantic version string (MUST follow semver).
  \item[\texttt{meta.description}] --- Human-readable summary of the Kit's purpose.
  \item[\texttt{compiler.executable}] --- Path (relative to Kit root) to the
    compiler binary.
\end{description}

\paragraph{Optional but recommended fields:}

\begin{description}
  \item[\texttt{compatibility.gert-core}] --- Semver range specifying which
    gert core versions this Kit is compatible with. If omitted, the Kit claims
    compatibility with all versions (not recommended).
  \item[\texttt{schema.apiVersion-suffix}] --- The suffix appended to
    \texttt{runbook/v2+} in runbooks authored with this Kit. Example:
    \texttt{compliance/v1} results in \texttt{apiVersion: runbook/v2+compliance/v1}.
  \item[\texttt{schema.json-schema}] --- Path to the Kit's JSON Schema overlay.
  \item[\texttt{step-types}] --- List of Kit-defined step types with metadata
    about what core types they lower to.
\end{description}

% ============================================================
\section{Kit Versioning Contract}
\label{sec:kit-versioning}
% ============================================================

Domain Kits follow semantic versioning (semver). The version in
\texttt{meta.version} has the following contract:

\paragraph{Major version bump (X.0.0):}
Breaking changes to the authoring schema. Examples:

\begin{itemize}
  \item Removing a step type
  \item Renaming a required field
  \item Changing field semantics in a non-backward-compatible way
  \item Changing the \texttt{apiVersion-suffix} (e.g., \texttt{compliance/v1}
    $\to$ \texttt{compliance/v2})
\end{itemize}

\paragraph{Minor version bump (1.X.0):}
Additive, backward-compatible changes. Examples:

\begin{itemize}
  \item Adding a new step type
  \item Adding optional fields to an existing step type
  \item Adding a new validator or projection
  \item Improving error messages without changing schema
\end{itemize}

\paragraph{Patch version bump (1.2.X):}
Bug fixes with no schema changes. Examples:

\begin{itemize}
  \item Fixing a lowering bug that produces incorrect core YAML
  \item Fixing a validator that incorrectly rejects valid input
  \item Fixing a projection that produces malformed output
\end{itemize}

\paragraph{Pinning in runbooks:}

Runbook authors can pin to a specific Kit version in the runbook metadata:

\begin{minted}{yaml}
apiVersion: runbook/v2+compliance/v1
meta:
  kit:
    name: gert.compliance
    version: "1.2.0"              # Exact version pin
  # OR:
  kit:
    name: gert.compliance
    version: "^1.2.0"             # Semver range (any 1.x >= 1.2.0)
\end{minted}

If the installed Kit version does not satisfy the pin, \texttt{gert validate}
rejects the runbook with a structured error.

% ============================================================
\section{Example Kit: \texttt{gert.compliance}}
\label{sec:example-kit}
% ============================================================

To make this concrete, here is the complete file tree for a minimal compliance
Kit:

\begin{verbatim}
gert.compliance/
├── gert-kit.yaml
├── bin/
│   └── gert-kit-compliance       # Go binary (compiled from cmd/compiler/main.go)
├── schemas/
│   └── compliance-v1.json        # JSON Schema overlay
├── test/
│   ├── fixtures/
│   │   └── basic-approval.yaml   # Fixture runbook using approval-request
│   └── expected/
│       └── basic-approval.lowered.yaml  # Expected lowered output
└── README.md
\end{verbatim}

The \texttt{gert-kit-compliance} binary is invoked by gert when processing a
runbook with \texttt{apiVersion: runbook/v2+compliance/v1}. It reads Kit YAML
on stdin, validates it against the Kit schema, and writes lowered
\texttt{runbook/v2} YAML to stdout.

In the next chapter, we will define the Kit schema that describes the
\texttt{approval-request} step type.
